% Options for packages loaded elsewhere
\PassOptionsToPackage{unicode}{hyperref}
\PassOptionsToPackage{hyphens}{url}
%
\documentclass[
  11pt,
]{article}
\usepackage{lmodern}
\usepackage{amssymb,amsmath}
\usepackage{ifxetex,ifluatex}
\ifnum 0\ifxetex 1\fi\ifluatex 1\fi=0 % if pdftex
  \usepackage[T1]{fontenc}
  \usepackage[utf8]{inputenc}
  \usepackage{textcomp} % provide euro and other symbols
\else % if luatex or xetex
  \usepackage{unicode-math}
  \defaultfontfeatures{Scale=MatchLowercase}
  \defaultfontfeatures[\rmfamily]{Ligatures=TeX,Scale=1}
  \setmainfont[]{Roboto Light}
\fi
% Use upquote if available, for straight quotes in verbatim environments
\IfFileExists{upquote.sty}{\usepackage{upquote}}{}
\IfFileExists{microtype.sty}{% use microtype if available
  \usepackage[]{microtype}
  \UseMicrotypeSet[protrusion]{basicmath} % disable protrusion for tt fonts
}{}
\makeatletter
\@ifundefined{KOMAClassName}{% if non-KOMA class
  \IfFileExists{parskip.sty}{%
    \usepackage{parskip}
  }{% else
    \setlength{\parindent}{0pt}
    \setlength{\parskip}{6pt plus 2pt minus 1pt}}
}{% if KOMA class
  \KOMAoptions{parskip=half}}
\makeatother
\usepackage{xcolor}
\IfFileExists{xurl.sty}{\usepackage{xurl}}{} % add URL line breaks if available
\IfFileExists{bookmark.sty}{\usepackage{bookmark}}{\usepackage{hyperref}}
\hypersetup{
  pdftitle={Informe del 04 de junio de 2021},
  hidelinks,
  pdfcreator={LaTeX via pandoc}}
\urlstyle{same} % disable monospaced font for URLs
\usepackage[margin=1in]{geometry}
\usepackage{graphicx,grffile}
\makeatletter
\def\maxwidth{\ifdim\Gin@nat@width>\linewidth\linewidth\else\Gin@nat@width\fi}
\def\maxheight{\ifdim\Gin@nat@height>\textheight\textheight\else\Gin@nat@height\fi}
\makeatother
% Scale images if necessary, so that they will not overflow the page
% margins by default, and it is still possible to overwrite the defaults
% using explicit options in \includegraphics[width, height, ...]{}
\setkeys{Gin}{width=\maxwidth,height=\maxheight,keepaspectratio}
% Set default figure placement to htbp
\makeatletter
\def\fps@figure{htbp}
\makeatother
\setlength{\emergencystretch}{3em} % prevent overfull lines
\providecommand{\tightlist}{%
  \setlength{\itemsep}{0pt}\setlength{\parskip}{0pt}}
\setcounter{secnumdepth}{-\maxdimen} % remove section numbering
\usepackage{titling}
\usepackage{fancyhdr}
\usepackage{lastpage}



\topmargin=-0.45in
\evensidemargin=-1in
\oddsidemargin=0in
\textwidth=6.5in
\textheight=9.0in
\headsep=0.25in

\linespread{1}
 \pagestyle{fancy}
 \fancyhead[LO,LE]{\thetitle}
 \fancyhead[RO,RE]{\includegraphics[width=0.5cm]{./DSLab_logo_2.png}}
 \fancyfoot[LO,LE]{DSLab, methaodos.org, Academia Joven de España \thedate}
 \cfoot{}
 \fancyfoot[RE,RO]{Página\ \thepage\ de\ \protect\pageref*{LastPage}}

\renewcommand\headrulewidth{0.4pt}
\renewcommand\footrulewidth{0.4pt}


\pretitle{%
  \begin{center}
  \LARGE
  \vspace{-2.5cm}
  \includegraphics[width=3.2cm]{./methaodos.png}
  \hspace{3cm}
  \includegraphics[width=3.5cm]{./dslab.png}
  \hspace{3cm}
  \includegraphics[width=3cm]{./aje.png}\\ 
\vspace{1cm}
\textbf{Suficiencia Sanitaria y COVID-19}\\
}
\posttitle{\end{center}}
\usepackage[font=small,skip=8pt]{caption}
\usepackage{float}
\floatplacement{figure}{H}
\renewcommand{\contentsname}{Índice}
\usepackage{booktabs}
\usepackage{longtable}
\usepackage{array}
\usepackage{multirow}
\usepackage{wrapfig}
\usepackage{float}
\usepackage{colortbl}
\usepackage{pdflscape}
\usepackage{tabu}
\usepackage{threeparttable}
\usepackage{threeparttablex}
\usepackage[normalem]{ulem}
\usepackage{makecell}
\usepackage{xcolor}

\title{Informe del 04 de junio de 2021}
\author{}
\date{\vspace{-2.5em}}

\begin{document}
\maketitle

\renewcommand{\figurename}{Figura}
\renewcommand{\tablename}{Tabla}

\tableofcontents

\vspace{.5cm}

\begin{center}\rule{0.5\linewidth}{0.5pt}\end{center}

\vspace{.5cm}

Para analizar el comportamiento de la pandemia provocada por el virus
SARS-CoV-2 y la enfermedad COVID-19 en España, se han diseñado dos
\textbf{Indicadores de Suficiencia Sanitaria (ISS):} uno acumulado
\textbf{(ISSa)} y otro semanal \textbf{(ISSs)}. Además, se lleva a cabo
una estimación del número de \textbf{casos activos} desde que se inició
la pandemia.

\hypertarget{indicador-de-suficiencia-sanitaria-acumulado-issa}{%
\subsection{Indicador de Suficiencia Sanitaria acumulado
(ISSa)}\label{indicador-de-suficiencia-sanitaria-acumulado-issa}}

El ISSa está definido en términos de porcentaje. Desde esa perspectiva,
valores altos del ISSa (cercanos a 100) indicarán que el sistema
sanitario está en disposición de absorber el flujo de diagnósticos
confirmados de coronavirus para que puedan ser tratados según los
protocolos establecidos. Por el contrario, valores bajos del ISSa
(cercanos a 0) sintomatizarán limitaciones que impedirán la correcta
gestión del flujo de casos diagnosticados. Tomando como referencia un
análisis de los valores del ISSa de la primera ola, se deben alcanzar
valores de este indicador por encima del 95\% para poder dar por
controlada la situación de pandemia desde un punto de vista de
suficiencia sanitaria.

Del ISSa se presentan dos estimaciones:

\begin{itemize}
\item
  Una primera estimación de su valor desde julio de 2020 hasta la
  actualidad (que se corresponde con las denominadas segunda y tercera
  olas). \textbf{El ISSa de España correspondiente a este periodo de la
  pandemia se mantiene por encima del valor del 95\%} (actualmente ISSa
  = 98.26\%, ver tabla 1) desde hace cuatro meses, con una tendencia
  creciente por encima de ese valor (ver figura 1, izquierda).
  \textbf{Este comportamiento del valor del indicador es coherente con
  la consistencia temporal de las medidas, mayoritariamente reactivas,
  que se están tomando a lo largo de esta etapa, y especialmente las
  tomadas desde el mes de enero hasta la actualidad.} Entendemos por
  \textbf{medidas reactivas} aquellas que se toman para dar respuesta a
  valores extremos de determinados indicadores (como por ejemplo la tasa
  de incidencia a 14 días). El denominador común de estas medidas es que
  se relajan una vez los valores de los indicadores dejan de ser
  extremos.
\item
  Una segunda estimación de su valor a lo largo de toda la pandemia, es
  decir, tomando para su cálculo el periodo completo que abarca desde el
  mes de marzo de 2020 hasta la actualidad (ver figura 1, derecha). Su
  valor supera también el 95\% (actualmente ISSa = 97.26\%, ver tabla 1)
  desde hace más de tres meses. Esta estimación permite ver el avance
  completo de la pandemia.
\end{itemize}

También se calcula el ISSa para cada comunidad autónoma. Se puede
comprobar en la figura 2 que \textbf{el ISSa del periodo julio
2020-junio 2021 (segunda y tercera olas) tiene también un comportamiento
creciente en todas las autonomías desde hace cuatro meses.} De hecho, en
todas las comunidades autónomas el indicador toma valores por encima del
95\%.

\begin{table}[!h]

\caption{\label{tab:tabla}ISS acumulado por CCAA, desde julio 2020 y pandemia.}
\centering
\fontsize{9}{11}\selectfont
\begin{tabular}[t]{l|r|r}
\hline
Comunidad Autónoma & ISSa (\%) desde julio 2020 & ISSa (\%) pandemia\\
\hline
Andalucia & 98.57 & 98.61\\
\hline
Aragón & 99.09 & 99.14\\
\hline
Asturias & 99.33 & 99.37\\
\hline
Baleares & 99.56 & 99.58\\
\hline
C. Valenciana & 99.70 & 99.71\\
\hline
Canarias & 98.48 & 98.54\\
\hline
Cantabria & 99.08 & 99.17\\
\hline
Castilla La Mancha & 99.20 & 99.29\\
\hline
Castilla y León & 99.21 & 99.28\\
\hline
Cataluña & 99.12 & 99.21\\
\hline
Ceuta & 99.77 & 99.78\\
\hline
Extremadura & 99.37 & 99.41\\
\hline
Galicia & 99.10 & 99.18\\
\hline
La Rioja & 98.98 & 99.13\\
\hline
Madrid & 99.03 & 99.13\\
\hline
Melilla & 99.09 & 99.10\\
\hline
Murcia & 99.40 & 99.42\\
\hline
Navarra & 99.27 & 99.34\\
\hline
País Vasco & 98.73 & 98.87\\
\hline
España & 98.26 & 97.26\\
\hline
\end{tabular}
\end{table}

\begin{figure}
\centering
\includegraphics{20210604_COVID19_DSLAB_informe_ejecutivo_files/figure-latex/fig_iss3-1.pdf}
\caption{\label{fig:fig_iss3} ISS acumulado, desde julio 2020 y
pandemia.}
\end{figure}

\vspace{0.2cm}

\begin{figure}
\centering
\includegraphics{20210604_COVID19_DSLAB_informe_ejecutivo_files/figure-latex/fig17a_res-1.pdf}
\caption{\label{fig:fig17a_res} ISS acumulado en las Comunidades
Autónomas, desde julio 2020.}
\end{figure}

\vspace{0.2cm}

\begin{figure}
\centering
\includegraphics{20210604_COVID19_DSLAB_informe_ejecutivo_files/figure-latex/fig17a_res2-1.pdf}
\caption{\label{fig:fig17a_res} ISS acumulado en las Comunidades
Autónomas, desde el inicio de la pandemia.}
\end{figure}

\newpage
\setcounter{page}{6}

\hypertarget{indicador-de-suficiencia-sanitaria-semanal-isss}{%
\subsection{Indicador de Suficiencia Sanitaria semanal
(ISSs)}\label{indicador-de-suficiencia-sanitaria-semanal-isss}}

El ISSs es un indicador que toma valores positivos, siendo clave superar
un valor de 1 para garantizar la suficiencia del sistema. Sucesivos
valores semanales por debajo de 1 implican posibles saturaciones del
sistema sanitario, mientras que si el indicador se sitúa por encima de 1
evidencia mejores niveles de suficiencia. Para que el sistema sanitario
esté en condiciones de estabilizarse es necesario que el ISSa mantenga
una tendencia creciente y que el ISSs se sitúe de manera sostenida por
encima de 1. \textbf{Desde una perspectiva nacional, el ISSs se ha
mantenido en los últimos dos meses por encima del valor de 1
(actualmente ISSs = 1.06, ver tabla 2), evidenciando esta situación una
clara mejora de la suficiencia sanitaria (ver figura 4).} Desagregando
por comunidades autónomas, todas ellas presentan un comportamiento
similar (ver figura 5). Estos valores continuados por encima de 1 son
coherentes con la tendencia creciente por encima del 95\% del indicador
acumulado ISSa.

\begin{table}[!h]

\caption{\label{tab:tabla2}ISS semanal por CCAA, desde julio 2020 y pandemia.}
\centering
\fontsize{9}{11}\selectfont
\begin{tabular}[t]{l|r|r}
\hline
Comunidad Autónoma & ISSs desde julio 2020 & ISSs pandemia\\
\hline
Andalucia & 0.95 & 0.95\\
\hline
Aragón & 1.18 & 1.18\\
\hline
Asturias & 1.07 & 1.07\\
\hline
Baleares & 0.98 & 0.98\\
\hline
C. Valenciana & 0.85 & 0.85\\
\hline
Canarias & 1.07 & 1.07\\
\hline
Cantabria & 0.87 & 0.87\\
\hline
Castilla La Mancha & 1.14 & 1.14\\
\hline
Castilla y León & 1.03 & 1.03\\
\hline
Cataluña & 1.21 & 1.21\\
\hline
Ceuta & 0.36 & 0.36\\
\hline
Extremadura & 1.00 & 1.00\\
\hline
Galicia & 1.00 & 1.00\\
\hline
La Rioja & 0.71 & 0.71\\
\hline
Madrid & 1.18 & 1.18\\
\hline
Melilla & 1.55 & 1.55\\
\hline
Murcia & 0.86 & 0.86\\
\hline
Navarra & 1.14 & 1.14\\
\hline
País Vasco & 1.04 & 1.04\\
\hline
España & 1.06 & 1.06\\
\hline
\end{tabular}
\end{table}

\begin{figure}
\centering
\includegraphics{20210604_COVID19_DSLAB_informe_ejecutivo_files/figure-latex/fig_iss2-1.pdf}
\caption{\label{fig:fig_iss2} ISS semanal, desde julio 2020 y pandemia.}
\end{figure}

\begin{figure}
\centering
\includegraphics{20210604_COVID19_DSLAB_informe_ejecutivo_files/figure-latex/figura_ISSSD_ccaa-1.pdf}
\caption{\label{fig:figura_ISSSD_ccaa} ISS semanal por Comunidades
Autónomas.}
\end{figure}

\clearpage

\setcounter{page}{9}

\hypertarget{estimaciuxf3n-de-casos-activos}{%
\subsection{Estimación de casos
activos}\label{estimaciuxf3n-de-casos-activos}}

Las \textbf{figuras 6 y 7} muestran tendencias que resumen de forma muy
ajustada la evolución de la pandemia en España. En concreto, presentan
\textbf{la estimación del número de casos diagnosticados que permanecen
activos cada semana}, tanto a nivel nacional como por comunidades
autónomas. \textbf{Todas las comunidades autónomas han vuelto a la senda
del decrecimiento de casos activos. Esta situación es coherente con el
mantenimiento en el tiempo de las medidas reactivas y con el efecto de
las medidas preventivas (como por ejemplo la vacunación).}

\vspace{0.2cm}

\begin{figure}
\centering
\includegraphics{20210604_COVID19_DSLAB_informe_ejecutivo_files/figure-latex/fig1-1.pdf}
\caption{\label{fig:fig17a_res} Estimación de casos activos en España
desde el inicio de la pandemia.}
\end{figure}

\begin{figure}
\centering
\includegraphics{20210604_COVID19_DSLAB_informe_ejecutivo_files/figure-latex/figura_activos_ccaa-1.pdf}
\caption{\label{fig:figura_activos_ccaa} Estimación de casos activos por
Comunidades Autónomas.}
\end{figure}

\newpage
\setcounter{page}{11}

\hypertarget{definiciones}{%
\subsection{Definiciones}\label{definiciones}}

Definiciones básicas:

\begin{itemize}
\item
  \(Confirmados_s\): número de diagnósticos confirmados acumulados a
  fecha \(s\).
\item
  \(NuevosCasos_{s} = Confirmados_{s} - Confirmados_{s-1}\): nuevos
  contagiados la semana \(s\).
\item
  \(Fallecidos_s\): número de fallecimientos acumulados a fecha \(s\).
\item
  \(FallecidosSemanales_s = Fallecidos_{s} - Fallecidos_{s-1}\): número
  de fallecimientos la semana \(s\).
\item
  \(Recuperados_s\): número de casos recuperados acumulados a fecha
  \(s\).
\item
  \(RecuperadosSemanales_{s} = Recuperados_{s} - Recuperados_{s-1}\):
  número de recuperados en la semana \(s\).
\item
  \(ActivosSemanales_s = Confirmados_s - Fallecidos_s - Recuperados_s\).
  Casos activos al final de la semana \(s\).
\item
  \(AtendidosSemanales_{s} = RecuperadosSemanales_s + FallecidosSemanales_s\):
  Pacientes atendidos la semana \(s\).
\item
  \(Atendidos_{s} = Recuperados_s + Fallecidos_s\): Pacientes atendidos
  acumulados a fecha \(s\).
\item
  \textbf{ISS acumulado a semana \(s\)}
  \[ISS\_Acumulado_s = \frac{Atendidos_s}{Confirmados_s} \times 100\]
  Este indicador mide el estrés acumulado del sistema. Cuanto más
  próximo a 100, más desestresado está el sistema.
\item
  \textbf{ISS Semanal, semana \(s\)}
\end{itemize}

\[
ISS_s = \frac{AtendidosSemanales_s}{NuevosCasos_{s}}
\]

\vspace{4cm}

\begin{center}
\includegraphics{cc.png}
\end{center}

\vspace{0.2cm}

\begin{center}
Esta obra está bajo una licencia de Creative Commons 
\end{center}

\begin{center}
Reconocimiento-CompartirIgual 4.0 Internacional.
\end{center}

\end{document}
